\section{Definitions}
\label{sec:definitions}

This section defines the list model itself. A list is either empty or a
single value paired with a smaller list, and every operation used later in
the article --- size, append, slicing, indexing, sum, and product --- is
defined by recursion on that same head/tail decomposition.

\subsection{List Construction}

Let $\mathbb{L}$ be the set of all lists over a set $S$. A list is either
the empty $L_{e}$ or a non-empty list $L_{node}$, as follows:

\subsection{Empty List}

Let's define an empty list $L_{e}$:

\begin{equation*}
\begin{aligned}
L_{e} & \in \mathbb{L} \\
L_{e} & = [] \\
\end{aligned}
\end{equation*}

\subsection{Recursive Definition of List}

A non-empty list packages one value, its \textbf{head}, together with the
remainder of the list, its \textbf{tail}, which is itself a smaller list.
Every list in $\mathbb{L}$ is either the single empty list or one of these
head/tail pairings, so the definition below builds the whole set
$\mathbb{L}$ out of the already-defined $L_e$ plus this one construction
rule:

\begin{equation*}
\begin{aligned}
&\text{head} & \in \mathbb{S} \\
&\text{tail} & \in \mathbb{L} \\
&L_{node}(\text{head}, \text{tail}) & \in \mathbb{L}_{node} \\
&\mathbb{L} = \{ L_e \} \cup \{ L_{node}(\text{head}, \text{tail}) & \mid \text{head} \in \mathbb{S},\ \text{tail} \in \mathbb{L} \} \\
\end{aligned}
\end{equation*}

\textbf{Termination and cyclic references.} Because all lists in this
model are immutable, each application of $L_{\text{node}}(\text{head},
\text{tail})$ produces a distinct structural value without the
possibility of cyclic references. Recursive functions over $\mathbb{L}$
terminate naturally, as a strictly decreasing structure defines size.

\subsection{Elements Access and Indexing}

The head/tail decomposition gives direct access to a list's first element
and its remaining sublist. Indexing extends this one step at a time:
position $0$ is the head, and position $n > 0$ is found by re-indexing
the tail at position $n - 1$, so reaching index $n$ costs $n$ recursive
tail steps. The last element is the value at the final valid index,
$|L| - 1$.

\begin{equation*}
\begin{aligned}
\text{if } L_{node} = [v_0, v_1, \dots, v_{n-1} ] & \implies L_{node} = (\text{head}: v_0, \text{tail}: [v_1, \dots, v_{n-1}]) \\
\operatorname{head}(L_{node}) & = v_0 \\
\operatorname{tail}(L_{node}) & = [v_1, \dots, v_{n-1}] \\
\operatorname{last}(L_{node}) & = L_{node(|L| - 1)} \\
L_{node(0)} & = L_{(0)} = \operatorname{head}(L_{node}) \\
L_{node(n)} & = L_{(n)} = \operatorname{tail}(L_{node})(n - 1) \quad \forall\, n > 0 \\
\end{aligned}
\end{equation*}

\subsection{List Size}

With the structure of lists defined, we now introduce a recursive
definition for their size (or length). We define the size of a list $L$,
$|L|$ as follows:

\begin{equation*}
|L| = \begin{cases}
0 & \text{if } L = L_{e} \\
1 + |\operatorname{tail}(L)| & \text{otherwise} \\
\end{cases}
\end{equation*}

The size of a list is zero for the empty list, or one plus the size of
its tail otherwise. Proved in the native stainless library in
\texttt{stainless.collection.List}.

\subsection{List Append}

Let $A, B \in \mathbb{L}$ over some set $S$. The append operation
$A \mathbin{\texttt{++}} B$ is defined recursively as:

\begin{equation*}
\begin{aligned}
A \mathbin{\texttt{++}} B =
\begin{cases}
B & \text{if } A = L_e \\
L_{node}(\operatorname{head}(A), \operatorname{tail}(A) \mathbin{\texttt{++}} B) & \text{otherwise}
\end{cases}
\end{aligned}
\end{equation*}

{\sloppy Appending $B$ onto an empty list yields $B$; appending it onto a
non-empty list keeps $A$'s head in place and appends $B$ onto $A$'s tail.
Proved in the native stainless library in \texttt{stainless.collection.List}.\par}

\subsection{List Slice}
\label{sec:list-slice}

Let $L = [v_0, v_1, \dots, v_{n-1}]$, $i, j \in \mathbb{N}$, with
$i \leq j < n$.

\begin{equation*}
L[i \dots j] := [ L_k \mid k \in \mathbb{N},\ i \leq k \leq j ]
\end{equation*}

The slice from $i$ to $j$ keeps exactly the elements at positions $i$
through $j$, in order. The implementation of \texttt{slice} is available
in
\href{https://github.com/thiagomata/prime-numbers/blob/list-article-v1.0.0/src/main/scala/v1/chapter3/list/ListUtils.scala#slice}{\texttt{ListUtils}}.
The full Scala verification code is in Appendix~A.3.

\subsection{List Sum}

Let $\operatorname{sum} : \mathbb{L} \implies \mathbb{S}$ be a recursively
defined function:

\begin{equation*}
\operatorname{sum}(L) =
\begin{cases} 0 & \text{if } L = L_e \\
\operatorname{head}(L) + \operatorname{sum}(\operatorname{tail}(L)) & \text{otherwise} \\
\end{cases}
\end{equation*}

The sum of an empty list is zero; the sum of a non-empty list is its head
plus the sum of its tail. The implementation of \texttt{sum} is available
in
\href{https://github.com/thiagomata/prime-numbers/blob/list-article-v1.0.0/src/main/scala/v1/chapter3/list/ListUtils.scala#sum}{\texttt{ListUtils}}.
The full Scala verification code is in Appendix~A.7.

\subsection{List Product}

Let $\operatorname{product} : \mathbb{L} \implies \mathbb{S}$ be a
recursively defined function:

\begin{equation*}
\operatorname{product}(L) =
\begin{cases} 1 & \text{if } L = L_e \\
\operatorname{head}(L) \cdot \operatorname{product}(\operatorname{tail}(L)) & \text{otherwise} \\
\end{cases}
\end{equation*}

The product of an empty list is one; a non-empty list's product is its
head times the product of its tail. The implementation of
\texttt{product} is available in
\href{https://github.com/thiagomata/prime-numbers/blob/list-article-v1.0.0/src/main/scala/v1/chapter3/list/properties/ListProduct.scala}{\texttt{ListProduct}}.
The full Scala verification code is in Appendices~A.11 through A.15.
